\section{Class functions, Adams operations, and finite-group extensions}
\label{sec:chebotarev}

The root-of-unity formulas of \Cref{sec:finite-part} describe the
abelian side of the primitive-orbit calculus. This section is the
non-abelian analogue of that root-of-unity decomposition. In the
abelian case, cyclic reconstruction is governed by ordinary Fourier
analysis on roots of unity. For finite-group extensions the correct
replacement is the Peter--Weyl decomposition of class functions. The
new feature is that primitive extraction is no longer governed by
M\"obius inversion alone: it is M\"obius inversion intertwined with the
Adams operations \(g\mapsto g^m\). This is the mechanism that turns
twisted determinant data into primitive Peter--Weyl series. At the
orbit-counting level, the Chebotarev asymptotics for finite homogeneous
extensions of shifts go back to Noorani--Parry
\cite{NooraniParry1992ChebotarevShifts}; for the character-ring
background and Adams operations we use the standard representation-ring
formalism of Serre \cite{Serre1977LinearRepresentations}.

\subsection{Class-function logarithms and primitive inversion}

For a conjugacy class \(C\subset G\), we write \(\chi_\varrho(C)\) for
the common value of \(\chi_\varrho\) on \(C\), and \(\mathbf 1_C\) for
the indicator function of \(C\). If \(x\in\operatorname{Fix}(\sigma^n)\)
corresponds to the closed edge path \(e_0\cdots e_{n-1}\), define
\[
  \tau_n(x):=\tau(e_0)\cdots\tau(e_{n-1})\in G.
\]
For a class function \(\phi:G\to\CC\), set
\begin{equation}\label{eq:class-logarithm-def}
  T_{\phi,n}:=\sum_{x\in\operatorname{Fix}(\sigma^n)}\phi(\tau_n(x)),
  \qquad
  \mathcal L_\phi(z):=\sum_{n\ge 1}\frac{T_{\phi,n}}{n}z^n,
\end{equation}
and
\begin{equation}\label{eq:primitive-class-series-def}
  \Pi_\phi(z):=\sum_{\gamma\in\cP}\phi(\tau(\gamma))z^{|\gamma|}.
\end{equation}
Since \(\phi\) is bounded and \(p_n(A)\ll \lambda^n/n\), the series
\(\Pi_\phi(z)\) converges absolutely for \(|z|<\lambda^{-1}\).

\begin{proposition}[Conjugacy-class expansion]
  \label{prop:class-indicator-expansion}
  For every conjugacy class \(C\subset G\) and every \(g\in G\),
  \[
    \mathbf 1_C(g)
    =\frac{|C|}{|G|}
      \sum_{\varrho\in\Irr(G)}
        \overline{\chi_\varrho(C)}\,\chi_\varrho(g)
    =\frac{|C|}{|G|}
      \sum_{\varrho\in\Irr(G)}
        \chi_\varrho(C^{-1})\,\chi_\varrho(g).
  \]
\end{proposition}

\begin{proof}
  The irreducible characters form an orthonormal basis of the class
  functions on \(G\) for the inner product
  \[
    \langle f,h\rangle
    :=\frac{1}{|G|}\sum_{g\in G}f(g)\overline{h(g)}.
  \]
  Hence
  \[
    \mathbf 1_C
    =\sum_{\varrho\in\Irr(G)}
      \langle \mathbf 1_C,\chi_\varrho\rangle \chi_\varrho.
  \]
  Since \(\chi_\varrho\) is constant on \(C\),
  \[
    \langle \mathbf 1_C,\chi_\varrho\rangle
    =\frac{1}{|G|}\sum_{g\in C}\overline{\chi_\varrho(g)}
    =\frac{|C|}{|G|}\overline{\chi_\varrho(C)}
    =\frac{|C|}{|G|}\chi_\varrho(C^{-1}),
  \]
  which gives the claim.
\end{proof}

\begin{proposition}[Peter--Weyl determinant factorisation]
  \label{prop:kernel-peter-weyl-block-diagonalisation}
  Let \(\widetilde A\) be the adjacency matrix of the skew-product
  extension attached to \(\tau\). Then
  \[
    \det(I-z\widetilde A)
    =
    \prod_{\varrho\in\Irr(G)}
      \det(I-zA_\varrho)^{\dim\varrho}.
  \]
\end{proposition}

\begin{proof}
  The regular representation decomposes as
  \[
    \CC[G]
    \cong
    \bigoplus_{\varrho\in\Irr(G)}
      V_\varrho^{\oplus \dim\varrho}.
  \]
  Under this decomposition, \(\widetilde A\) becomes block diagonal with
  blocks \(A_\varrho\otimes I_{\dim\varrho}\). Taking determinants gives
  the stated factorisation.
\end{proof}

\begin{proposition}[Twisted spectral-radius bound]
  \label{prop:twisted-spectral-radius-bound}
  For every irreducible unitary representation
  \(\sigma\in\Irr(G)\),
  \[
    \rho(A_\sigma)\le \rho(\widetilde A)=\lambda.
  \]
\end{proposition}

\begin{proof}
  Let \(\alpha\in\spec(A_\sigma)\). Then
  \(\det(I-\alpha^{-1}A_\sigma)=0\), so
  \Cref{prop:kernel-peter-weyl-block-diagonalisation} gives
  \[
    \det(I-\alpha^{-1}\widetilde A)
    =
    \prod_{\varrho\in\Irr(G)}
      \det(I-\alpha^{-1}A_\varrho)^{\dim\varrho}
    =0.
  \]
  Hence \(\alpha\in\spec(\widetilde A)\), and therefore
  \(\rho(A_\sigma)\le \rho(\widetilde A)\).

  To compute \(\rho(\widetilde A)\), observe that admissible
  length-\(n\) paths in the skew product are in bijection with pairs
  \((g,\gamma)\), where \(g\in G\) and \(\gamma\) is an admissible
  length-\(n\) path in the base graph: once \(g\) and \(\gamma\) are
  fixed, the lift is unique. Thus the number of admissible
  length-\(n\) paths in the extension is exactly \(|G|\) times the
  corresponding number in the base. The two shifts therefore have the
  same topological entropy. For a finite directed graph, that entropy is
  the logarithm of the spectral radius of its adjacency matrix, hence
  \(\rho(\widetilde A)=\rho(A)=\lambda\).
\end{proof}

\noindent
Up to this point, the material is representation-theoretic
infrastructure. The specific content of the present paper begins when
the determinant data are pushed to the primitive level by
Adams--M\"obius inversion.

\begin{theorem}[Class-function determinant linearisation]
  \label{thm:class-function-linearisation}
  Let \(\phi:G\to\CC\) be a class function and write
  \[
    \phi=\sum_{\varrho\in\Irr(G)}\widehat\phi(\varrho)\chi_\varrho
  \]
  as in \eqref{eq:class-function-fourier}. Then, for \(|z|<\lambda^{-1}\),
  \begin{equation}\label{eq:class-logarithm-determinant-linearisation}
    \mathcal L_\phi(z)
    =
    \sum_{\varrho\in\Irr(G)}
      \widehat\phi(\varrho)\,
      \log\det(I-zA_\varrho)^{-1}.
  \end{equation}
\end{theorem}

\begin{proof}
  For \(n\ge 1\), expanding \(\phi\) in the irreducible character basis
  gives
  \[
    T_{\phi,n}
    =\sum_{x\in\operatorname{Fix}(\sigma^n)}
      \sum_{\varrho\in\Irr(G)}
        \widehat\phi(\varrho)\chi_\varrho(\tau_n(x))
    =\sum_{\varrho\in\Irr(G)}
      \widehat\phi(\varrho)\Tr(A_\varrho^n),
  \]
  because \(\Tr(A_\varrho^n)\) is the sum of \(\chi_\varrho\) over all
  period-\(n\) points. Therefore,
  \[
    \mathcal L_\phi(z)
    =\sum_{n\ge 1}\frac{z^n}{n}
      \sum_{\varrho\in\Irr(G)}
        \widehat\phi(\varrho)\Tr(A_\varrho^n).
  \]
  Moreover,
  \[
    \sum_{n\ge 1}\frac{|z|^n}{n}\abs{T_{\phi,n}}
    \le
    \norm{\phi}_\infty
    \sum_{n\ge 1}\frac{\Tr(A^n)}{n}|z|^n
    =
    \norm{\phi}_\infty \log\det(I-|z|A)^{-1}
    <\infty
  \]
  for \(|z|<\lambda^{-1}\). Since
  \[
    \abs{T_{\phi,n}}
    \le \norm{\phi}_\infty \#\operatorname{Fix}(\sigma^n)
    =\norm{\phi}_\infty \Tr(A^n),
  \]
  the series is absolutely convergent for \(|z|<\lambda^{-1}\), and we
  may exchange the order of summation to obtain
  \[
    \mathcal L_\phi(z)
    =
    \sum_{\varrho\in\Irr(G)}
      \widehat\phi(\varrho)
      \sum_{n\ge 1}\frac{\Tr(A_\varrho^n)}{n}z^n
    =
    \sum_{\varrho\in\Irr(G)}
      \widehat\phi(\varrho)\,
      \log\det(I-zA_\varrho)^{-1}.
  \]
  This proves \eqref{eq:class-logarithm-determinant-linearisation}.
\end{proof}

\begin{corollary}[Adams--M\"obius primitive inversion]
  \label{cor:class-function-adams-mobius}
  For every class function \(\phi:G\to\CC\) and every \(|z|<\lambda^{-1}\),
  \begin{equation}\label{eq:class-adams-mobius-inversion}
    \Pi_\phi(z)
    =
    \sum_{m\ge 1}\frac{\mu(m)}{m}\,
      \mathcal L_{\psi^m\phi}(z^m).
  \end{equation}
\end{corollary}

\begin{proof}
  For \(\gamma\in\cP\) with length \(\ell=|\gamma|\), the \(r\)-fold
  repetition of \(\gamma\) contributes \(\ell\) points to
  \(\operatorname{Fix}(\sigma^{r\ell})\), and each such point has group
  label conjugate to \(\tau(\gamma)^r\). Since \(\phi\) is a class
  function,
  \[
    \mathcal L_\phi(z)
    =
    \sum_{\gamma\in\cP}\sum_{r\ge 1}
      \frac{\phi(\tau(\gamma)^r)}{r}z^{r|\gamma|}.
  \]
  Hence
  \begin{align*}
    \sum_{m\ge 1}\frac{\mu(m)}{m}\mathcal L_{\psi^m\phi}(z^m)
    &=
    \sum_{m\ge 1}\frac{\mu(m)}{m}
      \sum_{\gamma\in\cP}\sum_{r\ge 1}
        \frac{\phi(\tau(\gamma)^{mr})}{r}z^{mr|\gamma|} \\
    &=
    \sum_{\gamma\in\cP}\sum_{s\ge 1}
      \frac{\phi(\tau(\gamma)^s)}{s}z^{s|\gamma|}
      \sum_{m\mid s}\mu(m),
  \end{align*}
  where \(s=mr\). Since \(\sum_{m\mid s}\mu(m)=1\) for \(s=1\) and
  \(0\) otherwise, only the primitive term survives, yielding
  \[
    \sum_{m\ge 1}\frac{\mu(m)}{m}\mathcal L_{\psi^m\phi}(z^m)
    =
    \sum_{\gamma\in\cP}\phi(\tau(\gamma))z^{|\gamma|}
    =\Pi_\phi(z).
  \]
  This proves \eqref{eq:class-adams-mobius-inversion}.
\end{proof}

\subsection{Primitive Peter--Weyl series}

Applying \Cref{cor:class-function-adams-mobius} to irreducible
characters produces the primitive Peter--Weyl coefficients directly
from the twisted determinants.

For \(\varrho\in\Irr(G)\), define
\begin{equation}\label{eq:primitive-peter-weyl-series}
  \Pi_\varrho(z):=\Pi_{\chi_\varrho}(z)
  =\sum_{n\ge 1}\pi_{n,\varrho}(A)z^n,
  \qquad
  \pi_{n,\varrho}(A)
  :=\sum_{\substack{\gamma\in\cP\\ |\gamma|=n}}
    \chi_\varrho(\tau(\gamma)).
\end{equation}
For \(m\ge 1\), write
\begin{equation}\label{eq:adams-character-decomposition}
  \psi^m\chi_\varrho
  =\sum_{\sigma\in\Irr(G)}a_{\varrho,\sigma}^{(m)}\chi_\sigma
\end{equation}
for the Adams decomposition in the character ring.

\begin{lemma}\label{lem:adams-coefficients-bounded}
  For every \(\varrho,\sigma\in\Irr(G)\) and every \(m\ge 1\),
  \[
    \abs{a_{\varrho,\sigma}^{(m)}}\le \dim\varrho.
  \]
\end{lemma}

\begin{proof}
  Since the irreducible characters are orthonormal,
  \[
    a_{\varrho,\sigma}^{(m)}
    =\langle \psi^m\chi_\varrho,\chi_\sigma\rangle.
  \]
  Hence
  \[
    \abs{a_{\varrho,\sigma}^{(m)}}
    \le \norm{\psi^m\chi_\varrho}_2\norm{\chi_\sigma}_2
    =\norm{\psi^m\chi_\varrho}_2.
  \]
  Moreover,
  \[
    \norm{\psi^m\chi_\varrho}_2^2
    =\frac{1}{|G|}\sum_{g\in G}\abs{\chi_\varrho(g^m)}^2
    \le (\dim\varrho)^2,
  \]
  because \(\abs{\chi_\varrho(h)}\le \dim\varrho\) for all \(h\in G\).
  Taking the square root gives the claim.
\end{proof}

\begin{remark}[Periodicity of Adams coefficients]
  \label{rem:adams-coefficients-periodic}
  Let \(\exp(G)\) denote the exponent of \(G\). Since
  \(g^{m+\exp(G)}=g^m\) for every \(g\in G\), one has
  \[
    \psi^{m+\exp(G)}\chi_\varrho=\psi^m\chi_\varrho.
  \]
  Consequently
  \(a_{\varrho,\sigma}^{(m+\exp(G))}=a_{\varrho,\sigma}^{(m)}\) for all
  \(m\), so only finitely many Adams coefficient vectors occur. In the
  estimates below we use this finiteness only through the uniform bound
  of \Cref{lem:adams-coefficients-bounded}.
\end{remark}

\begin{corollary}[Irreducible-channel determinant formula]
  \label{cor:primitive-peter-weyl-determinant}
  For every \(\varrho\in\Irr(G)\) and every \(|z|<\lambda^{-1}\),
  \begin{equation}\label{eq:primitive-peter-weyl-determinant}
    \Pi_\varrho(z)
    =
    \sum_{m\ge 1}\frac{\mu(m)}{m}
      \sum_{\sigma\in\Irr(G)}
        a_{\varrho,\sigma}^{(m)}
        \log\det(I-z^mA_\sigma)^{-1}.
  \end{equation}
\end{corollary}

\begin{proof}
  Apply \Cref{cor:class-function-adams-mobius} with
  \(\phi=\chi_\varrho\). Then
  \[
    \Pi_\varrho(z)
    =
    \sum_{m\ge 1}\frac{\mu(m)}{m}
      \mathcal L_{\psi^m\chi_\varrho}(z^m).
  \]
  By \eqref{eq:adams-character-decomposition} and the determinant
  linearisation \eqref{eq:class-logarithm-determinant-linearisation},
  \[
    \mathcal L_{\psi^m\chi_\varrho}(z^m)
    =
    \sum_{\sigma\in\Irr(G)}
      a_{\varrho,\sigma}^{(m)}
      \log\det(I-z^mA_\sigma)^{-1}.
  \]
  Substituting this identity proves
  \eqref{eq:primitive-peter-weyl-determinant}.
\end{proof}

\begin{remark}[Abelian toy model: the cyclic lift]
  \label{rem:abelian-toy-model-cyclic}
  Let \(G=\ZZ/q\ZZ\), and let \(\chi_j(h)=\omega_q^{jh}\) be its
  irreducible characters. For the canonical cyclic-lift cocycle one has
  \(A_{\chi_j}=\omega_q^jA\), and therefore
  \[
    \mathcal L_{\chi_j}(z)=\log\det(I-\omega_q^j zA)^{-1}.
  \]
  The Adams--M\"obius inversion
  \eqref{eq:class-adams-mobius-inversion} then reduces to the ordinary
  Fourier decomposition on the \(q\)-th roots of unity. This is the
  concrete bridge between the abelian formulas of
  \Cref{sec:finite-part} and the present finite-group formalism.
\end{remark}

\begin{corollary}[Primitive Peter--Weyl trace formula]
  \label{cor:primitive-peter-weyl-trace}
  For every \(\varrho\in\Irr(G)\) and every \(n\ge 1\),
  \begin{equation}\label{eq:primitive-peter-weyl-trace}
    \pi_{n,\varrho}(A)
    =
    \frac{1}{n}\sum_{m\mid n}\mu(m)
      \sum_{\sigma\in\Irr(G)}
        a_{\varrho,\sigma}^{(m)}
        \Tr(A_\sigma^{n/m}).
  \end{equation}
\end{corollary}

\begin{proof}
  Extract the coefficient of \(z^n\) in
  \eqref{eq:primitive-peter-weyl-determinant}. Since
  \[
    [z^n]\log\det(I-z^mA_\sigma)^{-1}
    =
    \begin{cases}
      \Tr(A_\sigma^{n/m})/(n/m), & m\mid n, \\
      0, & m\nmid n,
    \end{cases}
  \]
  the stated formula follows immediately.
\end{proof}

\subsection{Conjugacy classes and explicit Mertens constants}

The primitive Peter--Weyl coefficients determine the conjugacy-class
counts by Fourier inversion on \(G\).

\begin{theorem}[Conjugacy-class primitive decomposition]
  \label{thm:class-primitive-decomposition}
  Let \(C\subset G\) be a conjugacy class. If \(p_{n,C}\) denotes the
  number of primitive orbits of length \(n\) whose \(G\)-label lies in
  \(C\), then for \(|z|<\lambda^{-1}\),
  \begin{equation}\label{eq:class-primitive-generating-function}
    \sum_{n\ge 1}p_{n,C}z^n
    =
    \frac{|C|}{|G|}
      \sum_{\varrho\in\Irr(G)}
        \chi_\varrho(C^{-1})\,\Pi_\varrho(z),
  \end{equation}
  and therefore
  \begin{equation}\label{eq:class-primitive-coefficients}
    p_{n,C}
    =
    \frac{|C|}{|G|}
      \sum_{\varrho\in\Irr(G)}
        \chi_\varrho(C^{-1})\,\pi_{n,\varrho}(A).
  \end{equation}
\end{theorem}

\begin{proof}
  By definition,
  \[
    \sum_{n\ge 1}p_{n,C}z^n
    =
    \sum_{\gamma\in\cP}\mathbf 1_C(\tau(\gamma))z^{|\gamma|}.
  \]
  Applying \Cref{prop:class-indicator-expansion} to the class function
  \(\mathbf 1_C\) and using the linearity of \(\Pi_\phi\) yields
  \eqref{eq:class-primitive-generating-function}. Comparing coefficients
  gives \eqref{eq:class-primitive-coefficients}.
\end{proof}

\begin{corollary}[Conjugacy-class Artin--M\"obius trace formula]
  \label{cor:class-artin-mobius-trace}
  For every conjugacy class \(C\subset G\) and every \(n\ge 1\),
  \begin{equation}\label{eq:class-artin-mobius-trace}
    p_{n,C}
    =
    \frac{|C|}{|G|\,n}
      \sum_{\varrho\in\Irr(G)}
        \chi_\varrho(C^{-1})
        \sum_{m\mid n}\mu(m)
          \sum_{\sigma\in\Irr(G)}
            a_{\varrho,\sigma}^{(m)}
            \Tr(A_\sigma^{n/m}).
  \end{equation}
\end{corollary}

\begin{proof}
  Substitute \eqref{eq:primitive-peter-weyl-trace} into
  \eqref{eq:class-primitive-coefficients}.
\end{proof}

\begin{theorem}[Explicit class Mertens theorem and Chebotarev density
  under twisted spectral gap]
  \label{thm:class-mertens-explicit}
  Assume that
  \[
    \eta:=\max_{\varrho\neq\mathbf 1}\rho(A_\varrho)<\lambda,
    \qquad
    \Lambda_*:=\max\{\Lambda(A),\eta,\lambda^{1/2}\}.
  \]
  Then, for every conjugacy class \(C\subset G\),
  \begin{equation}\label{eq:class-primitive-asymptotic}
    p_{n,C}
    =
    \frac{|C|}{|G|}\frac{\lambda^n}{n}
    +O\!\left(\frac{\Lambda_*^n}{n}\right).
  \end{equation}
  In particular, for all sufficiently large \(n\),
  \begin{equation}\label{eq:class-chebotarev-density}
    \frac{p_{n,C}}{p_n(A)}
    =
    \frac{|C|}{|G|}
    +O\!\left(\left(\frac{\Lambda_*}{\lambda}\right)^n\right).
  \end{equation}
  Moreover, the series \(\Pi_\varrho(\lambda^{-1})\) converges
  absolutely for every non-trivial \(\varrho\), and
  \begin{equation}\label{eq:class-mertens-harmonic}
    \sum_{n\le N}\frac{p_{n,C}}{\lambda^n}
    =
    \frac{|C|}{|G|}H_N
    +\mathcal M_C(A)-\frac{|C|}{|G|}\gamma
    +O\!\left(\frac{(\Lambda_*/\lambda)^N}{N}\right),
  \end{equation}
  where \(H_N=\sum_{n\le N}1/n\) and
  \begin{equation}\label{eq:class-mertens-constant}
    \mathcal M_C(A)
    =
    \frac{|C|}{|G|}\mathcal M(A)
    +\frac{|C|}{|G|}
      \sum_{\varrho\neq\mathbf 1}
        \chi_\varrho(C^{-1})\,\Pi_\varrho(\lambda^{-1}).
  \end{equation}
  In particular,
  \begin{equation}\label{eq:class-mertens-log}
    \sum_{n\le N}\frac{p_{n,C}}{\lambda^n}
    =
    \frac{|C|}{|G|}\log N+\mathcal M_C(A)+O(N^{-1}).
  \end{equation}
  Finally, for every non-trivial \(\varrho\),
  \begin{equation}\label{eq:class-mertens-determinant}
    \Pi_\varrho(\lambda^{-1})
    =
    \sum_{m\ge 1}\frac{\mu(m)}{m}
      \sum_{\sigma\in\Irr(G)}
        a_{\varrho,\sigma}^{(m)}
        \log\det(I-\lambda^{-m}A_\sigma)^{-1}.
  \end{equation}
  The implied constants in
  \eqref{eq:class-primitive-asymptotic} and
  \eqref{eq:class-mertens-harmonic} depend on \(A\), \(G\), and
  \(\tau\), but are uniform in the conjugacy class \(C\) and in \(n\)
  or \(N\), respectively.
\end{theorem}

\begin{proof}
  By \eqref{eq:class-primitive-coefficients},
  \[
    p_{n,C}
    =
    \frac{|C|}{|G|}p_n(A)
    +\frac{|C|}{|G|}
      \sum_{\varrho\neq\mathbf 1}
        \chi_\varrho(C^{-1})\,\pi_{n,\varrho}(A).
  \]
  For \(\varrho\neq\mathbf 1\), \Cref{cor:primitive-peter-weyl-trace}
  gives
  \[
    \pi_{n,\varrho}(A)
    =
    \frac{1}{n}\Tr(A_\varrho^n)
    +\frac{1}{n}
      \sum_{\substack{m\mid n\\ m>1}}\mu(m)
        \sum_{\sigma\in\Irr(G)}
          a_{\varrho,\sigma}^{(m)}\Tr(A_\sigma^{n/m}).
  \]
  The first term is \(O(\eta^n/n)\). For the second, by
  \Cref{lem:adams-coefficients-bounded} and the finiteness of
  \(\Irr(G)\),
  \[
    \sum_{\sigma\in\Irr(G)}
      a_{\varrho,\sigma}^{(m)}\Tr(A_\sigma^{n/m})
    =O(\lambda^{n/m}),
  \]
  uniformly in \(m\), because \(\rho(A_\sigma)\le \lambda\) for every
  \(\sigma\) (every eigenvalue of \(A_\sigma\) is an eigenvalue of the
  extension matrix \(\widetilde A\) by
  \Cref{prop:kernel-peter-weyl-block-diagonalisation}, and the finite
  skew-product extension has the same topological entropy as the base
  shift, so \(\rho(\widetilde A)=\lambda\)). Since \(n/m\le n/2\) for
  \(m>1\),
  \[
    \sum_{\substack{m\mid n\\ m>1}}\lambda^{n/m}
    \le \sum_{k=1}^{\lfloor n/2\rfloor}\lambda^k
    =O(\lambda^{n/2}).
  \]
  Hence
  \[
    \pi_{n,\varrho}(A)
    =O\!\left(\frac{\max\{\eta^n,\lambda^{n/2}\}}{n}\right).
  \]
  Combining this with \Cref{lem:primitive-orbit-asymptotic} proves
  \eqref{eq:class-primitive-asymptotic}.
  Since
  \[
    p_n(A)=\frac{\lambda^n}{n}+O\!\left(\frac{\Lambda_*^n}{n}\right),
  \]
  dividing \eqref{eq:class-primitive-asymptotic} by \(p_n(A)\) yields
  \eqref{eq:class-chebotarev-density} for all sufficiently large \(n\).

  Dividing \eqref{eq:class-primitive-asymptotic} by \(\lambda^n\) shows
  that \(\sum_{n\ge 1}\pi_{n,\varrho}(A)\lambda^{-n}\) converges
  absolutely for \(\varrho\neq\mathbf 1\), so
  \(\Pi_\varrho(\lambda^{-1})\) is well defined. Summing the identity
  for \(p_{n,C}\lambda^{-n}\) from \(1\) to \(N\), the trivial
  representation contributes
  \[
    \frac{|C|}{|G|}
      \sum_{n\le N}\frac{p_n(A)}{\lambda^n},
  \]
  while the non-trivial representations contribute
  \[
    \frac{|C|}{|G|}
      \sum_{\varrho\neq\mathbf 1}\chi_\varrho(C^{-1})
      \sum_{n\le N}\pi_{n,\varrho}(A)\lambda^{-n}.
  \]
  The tail of each non-trivial series is
  \(O((\Lambda_*/\lambda)^N/N)\), and
  \Cref{lem:primitive-orbit-asymptotic} also gives
  \[
    \sum_{n\le N}\frac{p_n(A)}{\lambda^n}
    =
    H_N+\sum_{n\ge 1}\left(\frac{p_n(A)}{\lambda^n}-\frac{1}{n}\right)
    +O\!\left(\frac{(\Lambda_*/\lambda)^N}{N}\right).
  \]
  By \Cref{cor:finite-part-mertens-asymptotic}, the convergent
  correction term equals \(\mathcal M(A)-\gamma\). This yields
  \eqref{eq:class-mertens-harmonic} and
  \eqref{eq:class-mertens-constant}. Since
  \(H_N=\log N+\gamma+O(N^{-1})\), equation
  \eqref{eq:class-mertens-log} follows. Finally,
  \eqref{eq:class-mertens-determinant} is
  \eqref{eq:primitive-peter-weyl-determinant} evaluated at
  \(z=\lambda^{-1}\).
\end{proof}

\begin{remark}[On the twisted spectral-gap hypothesis]
  \label{rem:twisted-spectral-gap-role}
  The condition
  \[
    \eta=\max_{\varrho\neq\mathbf 1}\rho(A_\varrho)<\lambda
  \]
  does not enter the determinant identities themselves: the formulas of
  \Cref{thm:class-function-linearisation,cor:class-function-adams-mobius,cor:primitive-peter-weyl-determinant}
  are valid throughout the disc \(|z|<\lambda^{-1}\) without it. The
  role of the hypothesis is exactly to permit evaluation at the Perron
  point \(z=\lambda^{-1}\) in the non-trivial channels and to convert
  the primitive Peter--Weyl series into explicit class Mertens
  constants. From the dynamical point of view this is the
  representation-by-representation analogue of the usual separation of
  the Perron eigenvalue from the rest of the spectrum. For structural
  information in this direction, compare the Chebotarev-type framework
  of \cite{ParryPollicott1986Chebotarev}, Parry's non-abelian
  Liv\v{s}ic theorem \cite{Parry1999LivsicNonAbelian}, and the analysis
  of Parry and Liv\v{s}ic obstructions in
  \cite{BoyleSchmieding2017FiniteGroupExtensions}. In this sense the
  hypothesis marks the boundary between formal primitive inversion and
  Perron-point asymptotics.
\end{remark}

\subsection{Non-abelian Fourier reconstruction}

The constants \(\mathcal M_C(A)\) therefore encode more than the leading
term of the class-counting asymptotic: they retain the full collection
of non-trivial twisted primitive values at the Perron point.

\begin{theorem}[Non-abelian Fourier reconstruction of class Mertens
  constants]
  \label{thm:class-mertens-fourier}
  Under the hypotheses of \Cref{thm:class-mertens-explicit}, define
  \[
    \Delta_C(A)
    :=\mathcal M_C(A)-\frac{|C|}{|G|}\mathcal M(A).
  \]
  Then
  \begin{equation}\label{eq:class-mertens-fourier-expansion}
    \Delta_C(A)
    =
    \frac{|C|}{|G|}
      \sum_{\varrho\neq\mathbf 1}
        \chi_\varrho(C^{-1})\,\Pi_\varrho(\lambda^{-1}),
  \end{equation}
  and for every non-trivial \(\sigma\in\Irr(G)\),
  \begin{equation}\label{eq:class-mertens-fourier-inversion}
    \Pi_\sigma(\lambda^{-1})
    =
    \sum_{C\subset G}\Delta_C(A)\chi_\sigma(C),
  \end{equation}
  where the sum runs over conjugacy classes. Moreover,
  \begin{equation}\label{eq:class-mertens-plancherel}
    \sum_{C\subset G}\frac{|G|}{|C|}\abs{\Delta_C(A)}^2
    =
    \sum_{\varrho\neq\mathbf 1}
      \abs{\Pi_\varrho(\lambda^{-1})}^2.
  \end{equation}
\end{theorem}

\begin{proof}
  Formula \eqref{eq:class-mertens-fourier-expansion} is simply
  \eqref{eq:class-mertens-constant} with the trivial term removed. To
  invert it, multiply by \(\chi_\sigma(C)\) and sum over conjugacy
  classes:
  \[
    \sum_{C\subset G}\Delta_C(A)\chi_\sigma(C)
    =
    \sum_{C\subset G}\frac{|C|}{|G|}
      \sum_{\varrho\neq\mathbf 1}
        \chi_\varrho(C^{-1})\Pi_\varrho(\lambda^{-1})\chi_\sigma(C).
  \]
  The orthogonality relation
  \[
    \sum_{C\subset G}\frac{|C|}{|G|}
      \chi_\varrho(C^{-1})\chi_\sigma(C)
    =\delta_{\varrho,\sigma}
  \]
  yields \eqref{eq:class-mertens-fourier-inversion}. Squaring
  \eqref{eq:class-mertens-fourier-expansion}, summing over conjugacy
  classes, and using the same orthogonality relation gives
  \eqref{eq:class-mertens-plancherel}.
\end{proof}

\begin{remark}
  \Cref{thm:class-mertens-fourier} is the non-abelian analogue of the
  root-of-unity reconstruction results from \Cref{sec:finite-part}. In
  the abelian case, one root-of-unity layer recovers the reduced
  spectral data by ordinary Fourier inversion. Here the full family of
  conjugacy-class Mertens constants recovers the non-trivial twisted
  primitive values \(\Pi_\varrho(\lambda^{-1})\) by non-abelian Fourier
  inversion.
\end{remark}

